%%%%%%%%%%%%%%%%%%%%%%%%%%%%%%%%%%%%%%%%%%%%%%%%%%%%%%%%%%%%%%%%%%%%%%%%%%%
%
% Generic template for TFC/TFM/TFG/Tesis
%
% $Id: resumen.tex,v 1.9 2015/06/05 00:10:31 macias Exp $
%
% By:
%  + Javier Macías-Guarasa. 
%    Departamento de Electrónica
%    Universidad de Alcalá
%  + Roberto Barra-Chicote. 
%    Departamento de Ingeniería Electrónica
%    Universidad Politécnica de Madrid   
% 
% Based on original sources by Roberto Barra, Manuel Ocaña, Jesús Nuevo,
% Pedro Revenga, Fernando Herránz and Noelia Hernández. Thanks a lot to
% all of them, and to the many anonymous contributors found (thanks to
% google) that provided help in setting all this up.
%
% See also the additionalContributors.txt file to check the name of
% additional contributors to this work.
%
% If you think you can add pieces of relevant/useful examples,
% improvements, please contact us at (macias@depeca.uah.es)
%
% You can freely use this template and please contribute with
% comments or suggestions!!!
%
%%%%%%%%%%%%%%%%%%%%%%%%%%%%%%%%%%%%%%%%%%%%%%%%%%%%%%%%%%%%%%%%%%%%%%%%%%%

\chapter*{Resumen}
\label{cha:resumen}
\markboth{Resumen}{Resumen}

\addcontentsline{toc}{chapter}{Resumen}

En este trabajo de Fin de Grado se ha explorado la capacidad que tienen diferentes modelos de aprendizaje automático no supervisado para separar señales temporales provocadas por 
fuentes de ruido de las generadas por fotones. El objetivo es identificar un modelo que supere la capacidad de clasificación de la Red Neuronal Convolucional (CNN) utilizada en el artículo que sirvió de punto de partida para este trabajo.
Para abordar esta tarea se han entrenado cuatro modelos diferentes con ambos tipos de señales (capturadas por el detector TES). Inicialmente, se ha analizado la naturaleza de las señales para luego aplicar las 
técnicas de preprocesado correspondientes. Una vez entrenados los modelos, se han evaluado bajo las mismas métricas, con el propósito de poder compararlos y determinar el modelo con mejor rendimiento discriminativo de ruido.

\textbf{Palabras clave:} \myThesisKeywords.

%%% Local Variables:
%%% TeX-master: "../book"
%%% End:


